\chapter{Heart Disease}

\section*{Patient focus}
Heart imaging helps doctors see how well your heart is pumping, whether the heart muscle is damaged, if the valves are working properly, and if the coronary arteries (the heart's own blood supply) are blocked. Different scans answer different questions: an echocardiogram uses sound waves to see the heart beating in real time, a CT scan can show blocked arteries, and an MRI can reveal scar tissue or inflammation in the heart muscle. The right test depends on what your doctor needs to find out.


\label{ch:ch61_heart_disease}

\section{Imaging Role}
Imaging evaluates ischemia, infarction, cardiomyopathy, valvular disease, heart failure, congenital disease, and vascular complications. Echocardiography, CT, MRI, nuclear imaging, and angiography should be selected for the specific clinical question.

\section{Choosing the Right Modality}

\begin{table}[h]
\centering
\caption{Cardiac imaging modality selection}
\begin{tabular}{|l|l|l|}
\hline
\textbf{Clinical question} & \textbf{Best modality} & \textbf{Key information} \\
\hline
Cardiomyopathy & CMR & Tissue characterisation, EF \\
Coronary anatomy & CT coronary angiography & Stenosis, calcium score \\
Ischemia vs scar & CMR (LGE) & Viability, infarction pattern \\
Myocarditis & CMR (T2, LGE, EGE) & Edema, inflammation, fibrosis \\
Nuclear perfusion & SPECT/PET & Ischemia, viability \\
Pericardial disease & Echo, CMR, CT & Effusion, constriction \\
Pre-op bypass planning & CT coronary angiography & Anatomy, graft suitability \\
Valvular disease & Echo (TTE/TEE) & Severity, mechanism, plan \\
Wall motion / EF & Echo (TTE) & Function, valves, pericardium \\
\hline
\end{tabular}
\end{table}

\section{Coronary Artery Disease}

\subsection{Coronary CT Angiography (CTA)}
\begin{itemize}
  \item \textbf{Technique}: ECG-gated, IV contrast, beta-blockade to achieve heart rate $<$65 bpm
  \item \textbf{Calcium scoring}: Agatston score---0 (no calcium), 1--99 (mild), 100--399 (moderate), $\geq$400 (severe)
  \item \textbf{Stenosis assessment}: high NPV for excluding significant disease (score 0 nearly excludes flow-limiting disease)
  \item \textbf{Limitations}: heavy calcification, stents, high/irregular heart rate, BMI
\end{itemize}

\subsection{Nuclear Myocardial Perfusion Imaging}
\begin{itemize}
  \item SPECT (Tc-99m sestamibi/tetrofosmin) at rest and stress
  \item PET (Rb-82 or N-13 ammonia) for quantitative perfusion and viability
  \item Quantify ischemic burden: $>$10\% LV ischemia generally warrants revascularization
\end{itemize}

\section{Cardiomyopathies}

\subsection{Dilated Cardiomyopathy (DCM)}
\begin{itemize}
  \item \textbf{Imaging}: enlarged LV (EDV $>$97\% percentile), reduced EF ($<$50\%), spherical remodelling
  \item \textbf{CMR LGE}: mid-wall streak (fibrosis pattern)---present in $\sim$30\% of non-ischemic DCM; confers higher arrhythmic risk
  \item \textbf{Differential}: ischemic cardiomyopathy (subendocardial or transmural LGE in coronary territory)
\end{itemize}

\subsection{Hypertrophic Cardiomyopathy (HCM)}
\begin{itemize}
  \item \textbf{Imaging}: asymmetric septal hypertrophy ($>$15 mm), LVOT obstruction (gradient $\geq$30 mmHg at rest or provocation)
  \item \textbf{CMR}: identifies apical HCM (missed on echo), fibrosis (LGE present in $\sim$60\%), and crypts
  \item \textbf{Sudden cardiac death risk}: extensive LGE ($\geq$15\% of LV mass) is a risk factor
\end{itemize}

\subsection{Restrictive Cardiomyopathy}
\begin{itemize}
  \item \textbf{Imaging}: small normal-sized chambers, biatrial enlargement, thickened ventricular walls, impaired diastolic filling
  \item \textbf{CMR}: helps identify cause---amyloid (diffuse subendocardial LGE, abnormal gadolinium kinetics), sarcoid (patchy mid-wall LGE), haemochromatosis (diffuse low T1 from iron)
\end{itemize}

\section{Myocarditis}
\begin{itemize}
  \item \textbf{CMR (Lake Louise criteria)}:
    \begin{enumerate}
      \item T2-weighted imaging: myocardial edema (high signal)
      \item Early gadolinium enhancement: hyperaemia
      \item Late gadolinium enhancement (LGE): typically epicardial or mid-wall in non-coronary distribution (inferolateral wall most common)
    \end{enumerate}
  \item \textbf{LGE pattern in myocarditis}: subepicardial, non-coronary distribution---distinct from ischemic (subendocardial/transmural, coronary territory)
\end{itemize}

\begin{pearlbox}[LGE patterns by aetiology]
\begin{enumerate}
  \item \textbf{Ischemic} (infarction): subendocardial or transmural, within coronary territory
  \item \textbf{DCM}: mid-wall septal streak
  \item \textbf{HCM}: patchy at RV insertion points or hypertrophied segments
  \item \textbf{Myocarditis}: subepicardial, non-coronary (inferolateral)
  \item \textbf{Cardiac amyloid}: diffuse subendocardial with blood pool enhancement
  \item \textbf{Cardiac sarcoid}: mid-wall, often septal and anterolateral
\end{enumerate}
\end{pearlbox}

\section{Pericardial Disease}
\begin{itemize}
  \item \textbf{Pericardial effusion}: echo is first-line; CT and MRI quantify and characterize. Tamponade requires urgent drainage
  \item \textbf{Constrictive pericarditis}: thickened pericardium ($>$4 mm), biatrial enlargement, septal bounce on cine MRI, respiratory variation in flow
  \item \textbf{Pericardial enhancement on CMR}: active inflammation---may respond to anti-inflammatory therapy rather than surgery
\end{itemize}

\section{Valvular Disease}
\begin{itemize}
  \item \textbf{Echo (TTE/TEE)}: first-line for valve assessment---severity, mechanism, haemodynamic consequences
  \item \textbf{CT}: valve calcification quantification (TAVR planning), annular sizing
  \item \textbf{CMR}: accurate regurgitant volume and fraction, especially for pulmonic and tricuspid valves where echo is less reliable
\end{itemize}

\section{Cardiac CT Protocol}
\begin{enumerate}
  \item \textbf{Pre-scan}: heart rate assessment, beta-blocker if $>$65 bpm, sublingual GTN
  \item \textbf{Calcium score}: non-contrast ECG-gated
  \item \textbf{CTA}: IV contrast (60--80 mL), bolus tracking to ascending aorta, ECG-gated (prospective or retrospective)
  \item \textbf{Post-processing}: MPR, MIP, VR for coronary assessment; CT-FFR if needed
\end{enumerate}

\section{Interpretation}
Assess function, wall motion, perfusion, late enhancement, chamber size, valves, pericardium, coronary anatomy, and extracardiac findings. Quantitative measurements require appropriate loading conditions and technique.

\section{Safety}
Consider rhythm, hemodynamic stability, contrast, renal function, implants, and radiation. Acute coronary syndrome, tamponade, dissection, severe outflow obstruction, or major pulmonary hypertension requires urgent clinical coordination.

\begin{pearlbox}[Cardiac imaging safety]
CMR safety depends on the exact implant or foreign body and its MR conditions; many cardiac devices are MR-conditional under specified protocols. Verify the device and follow the facility's MR safety process. For gadolinium, assess the specific agent and clinical indication: a clinically indicated ACR Group II agent is generally not withheld solely because eGFR is below 30, while Group I agents are avoided in severe renal impairment.
\end{pearlbox}

\clearpage
